\documentclass[12pt,a4paper]{article}
\usepackage[UTF8]{ctex}
\usepackage{geometry}
\usepackage{titlesec}
\usepackage{setspace}
\usepackage{indentfirst}
\usepackage[hidelinks]{hyperref}

\geometry{left=2.5cm,right=2.5cm,top=2.5cm,bottom=2.5cm}

\titleformat{\section}{\Large\bfseries}{}{0em}{}
\titleformat{\subsection}{\large\bfseries}{}{0em}{}

\setlength{\parindent}{2em}
\setlength{\parskip}{0.5em}

\title{\textbf{2026年度工作计划}}
\author{许达 \quad 未来院三室}
\date{}

\begin{document}

\maketitle

\section*{四、2026年度量化指标}

2026年度工作目标将坚持"可量化、可验收、可复盘"的原则设置指标体系，以确保工作推进有节奏、有抓手。

\subsection*{AI 赋能科研}
AI 赋能科研方面，全年组织 AI 培训或实操工作坊不少于 12 次。推动团队项目在周期或工时上实现不低于 20\% 的效率提升，并按季度形成能力评估报告。

\subsection*{Alphaqubit 项目}
Alphaqubit 方面，\textbf{牵头完成码距为3的实时AI量子纠错译码模型，满足科委项目的考核指标，支持FPGA和GPU两种部署方式}。完成迭代版本发布不少于 2 次，核心模块自动化测试覆盖率达到 60\% 以上，输出完整技术文档与使用手册，至少支撑或赋能团队项目 1 项，同时形成至少 1 次内部演示或跨团队技术交流。

\subsection*{量子大模型}
量子大模型方面，完成量子大模型原型平台 1 套，具备可演示、可部署能力；\textbf{牵头完成AI赋能量子编程软件原型，一次性成功率达到 50\%}；构建量子知识库文献条目不少于 5000 条；完成量子推理评测集不少于 200 题并形成评估报告；形成至少 1 个内部示范应用场景。

\subsection*{科研智能体}
科研智能体方面，形成科研智能体系统 1 套并达到可运行状态；至少落地 2 个科研场景应用；形成至少 2 次效果评估报告。

\subsection*{强化学习量子控制软件}
强化学习量子控制软件方面，形成可运行原型 1 套，完成不少于 3 种基线对比实验并输出 1 份报告，沉淀可复用软件模块不少于 2 个。

\subsection*{专利与论文专著}
专利与论文专著方面，完成量子加密算法专利申请不少于 2 项，完成论文投稿或预印本不少于 2 篇，完成专著或书稿/章节不少于 1 项，输出研究报告或综述不少于 2 份。

\end{document}
